\section{\bench{}}
\label{sec:benchmark}

\bench{} separates mechanism validation from photorealistic evaluation and physical deployment.
This design prevents perception errors from obscuring whether prefix-aware planning is correct, while retaining the complexity required to test a deployable embodied stack.

\subsection{Mini-\taskname{}: Exact Mechanism Tests}

Mini-\taskname{} comprises partially observed room graphs and small grid worlds with deterministic or stochastic sensors.
Tasks, arrival times, action prefixes, and evidence certificates can be enumerated exactly.
Each instance is constructed from one of four mechanism families:
(i) matched terminal evidence with different acquisition times;
(ii) long actions with high terminal gain versus short actions with early gain;
(iii) task-prior shifts that reverse the preferred evidence; and
(iv) macro-actions with controlled duration variance.
Dynamic programming provides an exact arrival-aware oracle.
The layer measures policy regret, evidence-time likelihood, and the approximation error induced by time binning or limited \jet{} rank.
It is not used to claim visual realism.

\subsection{Cross-Environment Simulation Suite}
\label{sec:cross-environment}

The simulation suite deliberately varies scene source, rendering fidelity, semantic diversity, and downstream goal type.
It includes large-scale procedural houses from ProcTHOR, scanned real-world interiors from HM3D, high-quality authored synthetic interiors from HSSD, high-fidelity Gaussian-splatting environments from Habitat-GS, and evolving HSSD histories from P4D-Bench
\citep{savva2019habitat,ramakrishnan2021hm3d,deitke2022procthor,khanna2024hssd,xia2026habitatgs}.
We use HM3D-OVON for open-vocabulary navigation, the active-exploration portion of OpenEQA for natural embodied questions, HSSD ObjectNav for object-search transfer, Habitat-GS PointNav for photorealistic renderer transfer, and P4D-Bench for temporally changing evidence
\citep{yokoyama2024hm3dovon,majumdar2024openeqa}.
ProcTHOR supports large-scale training and held-out compositions of the four \taskname{} task families.
These tasks are not pooled into a single leaderboard score; each retains its native goal representation and success rule where applicable.

\begin{table*}[t]
  \centering
  \scriptsize
  \setlength{\tabcolsep}{3pt}
  \resizebox{\linewidth}{!}{%
  \begin{tabular}{@{}lllll@{}}
    \toprule
    \textbf{Evaluation layer} & \textbf{Scene source} & \textbf{Representation} & \textbf{Downstream protocol} & \textbf{Primary role} \\
    \midrule
    Exact mechanism & Mini-\taskname{} & Room graph / grid & Enumerated service tasks & Oracle regret and timing \\
    Procedural simulation & ProcTHOR & Interactive houses & \taskname{} service catalog & Scale and composition \\
    Scanned simulation & HM3D & Textured mesh & HM3D-OVON & Open-vocabulary navigation \\
    Scanned / recorded & HM3D + real scans & RGB-D episodes & OpenEQA & Attribute and spatial QA \\
    Authored simulation & HSSD & Semantic 3D scenes & ObjectNav & Object-search transfer \\
    Photorealistic simulation & Habitat-GS & 3D Gaussian splats & PointNav & Renderer and control transfer \\
    Evolving simulation & HSSD & Versioned dynamic scenes & P4D-Bench & Temporal and stale-evidence stress \\
    Physical deployment & Indoor / outdoor & Three robot platforms & Four service families & Embodiment transfer \\
    \bottomrule
  \end{tabular}%
  }
  \caption{\textbf{Environment and benchmark coverage.}
  Every simulated protocol is wrapped by the same task-hidden arrival and atomic-interruption semantics; comparisons are made within, not across, heterogeneous benchmarks.}
  \label{tab:evaluation-matrix}
\end{table*}

In the \taskname{} wrapper, the benchmark-native goal is sampled from the declared service catalog but remains hidden during pre-task exploration.
When it arrives, the current macro-action is interrupted and the native downstream executor starts from the committed memory and interrupted pose.
OpenEQA questions are filtered to those answerable from embodied visual evidence and mapped to the typed attribute, spatial, and verification programs; questions that require undeclared external knowledge are excluded.
For P4D-Bench, the task catalog is sampled from a versioned HSSD scene history so that evidence can also become stale between observation sessions; this is reported separately from the within-action interruption analysis.

The agent receives only the observation and egomotion channels supported by the matched domain, standardized to RGB-D and a noisy pose belief where available.
It does not receive privileged global pose, semantic maps, navigation meshes, shortest-path distances, future frames, or unobserved object annotations.
Within each domain, open-vocabulary perception, local control, \rfm{} capacity, validators, candidate generation, arrival draws, and frozen post-arrival executors are identical across methods.
Training may use simulator cloning for counterfactual supervision; evaluation may not.

Scene splits are fixed before hyperparameter selection.
Training and validation scenes support model fitting and evidence-threshold calibration, while all headline results use unseen scenes or the official unseen split of the corresponding benchmark.
Every comparison is paired by domain, scene, start pose, service catalog, task seed, arrival seed, and sensor-noise seed.

\subsection{Physical Robot Deployment}
\label{sec:realworld-benchmark}

We deploy the complete \method{} stack on three commercial DEEP Robotics platforms: the wheel-legged LYNX M20 and the quadrupedal X30 and Lite3.
M20 is the primary quantitative platform; X30 and the LiDAR-equipped Lite3 evaluate cross-platform transfer across indoor and outdoor scenes.
The memory, \jet{} predictor, programs, and planner remain shared; only perception and low-level navigation adapters differ.
Hardware and physical interruption details appear in \cref{app:realworld}.

\begin{figure*}[t]
  \centering
  \placeholderfigure{4.3cm}{Benchmark overview: Mini-\taskname{}; ProcTHOR houses; HM3D/HM3D-OVON and OpenEQA; HSSD ObjectNav; Habitat-GS PointNav; dynamic HSSD/P4D-Bench; indoor/outdoor deployment on M20, X30, and Lite3.}
  \caption{\textbf{\bench{} design.}
  The controlled layer supplies exact oracles and diagnostic counterexamples.
  Procedural, scanned, authored-synthetic, Gaussian-splatting, evolving-scene, and physical deployments evaluate the same interruption semantics with noisy embodied perception and navigation.}
  \label{fig:benchmark}
\end{figure*}

\subsection{Task Families and Compositional Splits}

All natural-language requests compile to typed requirement programs.
The benchmark contains four families:
\begin{enumerate}[leftmargin=*,itemsep=2pt,topsep=3pt]
  \item \textbf{Category/instance navigation:} reach and confirm a named object, requiring a grounded location and reachability evidence.
  \item \textbf{Attribute QA:} answer a visible property such as color, material, state, or coarse shape from an identity-qualified, sufficiently resolved view.
  \item \textbf{Spatial QA:} answer a relation such as inside, on, near, or left-of from compatible entity bindings and geometrically consistent evidence.
  \item \textbf{Presence/absence verification:} output \textsc{Present}, \textsc{Absent}, or \textsc{Insufficient}. Absence requires a coverage certificate and a calibrated category detector; unsupported guessing is penalized.
\end{enumerate}

Splits are defined over requirement structure, not only language paraphrases.
Training includes single primitives and a subset of binary combinations.
The compositional test split holds out attribute--category, room--relation, and conjunction patterns while retaining the constituent primitives.
A separate parser track maps natural language to programs; the primary algorithmic comparison uses gold template programs so that parser errors do not confound exploration.

\subsection{Arrival Processes and Atomic Evaluation}

The benchmark defines five arrival regimes over the same real-time horizon:
\textbf{early-heavy}, with high initial hazard;
\textbf{memoryless}, with constant hazard;
\textbf{late-heavy}, with increasing hazard;
\textbf{burst/windowed}, with service peaks; and
\textbf{misspecified}, where the evaluation hazard differs from training.
Multiple arrival times can be resampled along a stored trajectory to reduce evaluation variance only when the policy itself was conditioned on the corresponding declared hazard.
A fixed policy cannot be retrospectively reweighted and presented as adaptation.

At each low-level step the evaluator applies the order
observation $\rightarrow$ memory commit $\rightarrow$ arrival check.
If interrupted, the action does not complete and future observations are inaccessible.
The evaluator snapshots the memory hash, pose belief, action prefix, task identifier, and downstream executor seed to make leakage audits possible.

\subsection{Utility and Metrics}

For each task family, raw post-arrival performance is mapped to $[0,1]$ using a predeclared normalization.
Navigation utility combines success within budget and normalized service cost; QA uses correctness with an abstention-aware score; absence verification uses selective risk; and recall-only utility measures performance with $B=0$.
The three primary metrics are:
\begin{align}
\mathrm{Recall}\text{-}\awu &=
\mathbb{E}_{z,\tau}[U_z(M_\tau,x_\tau;0)],\\
\mathrm{Service}\text{-}\awu &=
\mathbb{E}_{z,\tau}[U_z(M_\tau,x_\tau;B)],\\
\mathrm{WorstFamilyService}\text{-}\awu &=
\min_{g\in\mathcal{G}}\mathbb{E}[U_z\mid z\in g].
\end{align}
Here AWU denotes arrival-weighted utility.
Secondary metrics include the readiness curve and its area, task-family raw metrics, post-arrival success/cost/SPL, evidence grounding precision, absence risk--coverage, ECE/Brier score, blockage rate, memory size, latency, evidence-time C-index, and hitting-time calibration.
Coverage is reported only as a diagnostic, never as success.
